% ===========================================================================
%  AdaFoB -- Implementation Details / Experimental Setup
%  Requires: \usepackage{xcolor}
%            \newcommand{\needsnum}[1]{\textcolor{red}{[NUMBER PENDING: #1]}}
%            \newcommand{\needsrun}[1]{\textcolor{red}{[EXPERIMENT PENDING: #1]}}
%
%  FACTUAL WARNING (read before editing):
%    Abd-CT, SABS and BTCV are THREE NAMES FOR ONE DATASET (Synapse
%    syn3193805, Multi-Atlas Labeling Beyond the Cranial Vault). Do not write
%    that this work "replaced SABS with Abd-CT" -- that sentence is false and
%    a reviewer familiar with the ALPNet/FoB lineage will catch it instantly.
%    What changed is the PREPROCESSING PIPELINE, and that is what we declare.
% ===========================================================================

\section{Experimental Setup and Implementation Details}
\label{sec:implementation}

\subsection{Dataset}
\label{subsec:dataset}

We evaluate on \textbf{Abd-CT}, the abdominal CT benchmark used by the FSMIS
literature, comprising $30$ volumetric contrast-enhanced scans from the
Multi-Atlas Labeling Beyond the Cranial Vault challenge (Synapse
\texttt{syn3193805}). This resource appears in prior work under three
interchangeable names -- Abd-CT, SABS and BTCV -- and we use \emph{Abd-CT}
throughout for consistency with FoB~\cite{bo2026fob} and
SSL-ALPNet~\cite{ouyang2020ssl}. Following that lineage we evaluate the four
labels common to the abdominal benchmarks: liver, spleen, and left and right
kidney. We adopt the standard two protocols, Setting~I (test classes may occur
unlabelled in training backgrounds) and Setting~II (test classes are entirely
unseen), with $5$-fold cross-validation.

Abd-CT is an appropriate testbed for our hypothesis for two reasons beyond
comparability. First, the four evaluated organs span roughly an order of
magnitude in cross-sectional area, and within a single organ the axial area
varies strongly from the first to the last annotated slice; this within-benchmark
scale variance is precisely the axis that a fixed prompt budget cannot exploit.
Second, the label set contains structurally irregular organs (notably the
pancreas and the great vessels) that can be added to stratify results by shape
complexity without introducing a second imaging domain.

\paragraph{Preprocessing and a declared protocol deviation.}
FoB consumes volumes preprocessed by the pipeline of Ouyang
et~al.~\cite{ouyang2020ssl}, distributed as \texttt{sabs\_CT\_normalized}. Due
to distribution and formatting incompatibilities in our compute environment we
were unable to use that artefact and instead preprocess the source volumes
ourselves. Our pipeline clips Hounsfield units to $[-125, 275]$, rescales
linearly to $[0,255]$, resamples slices to $256\times256$ (bilinear for images,
nearest-neighbour for labels), and applies $z$-score normalisation using
statistics pooled across the corpus rather than per volume. Pooled statistics
are used because our volumes are cropped along the $z$-axis to the annotated
extent, which destabilises per-volume statistics.

We state explicitly that this is a \emph{deviation} from the reference
protocol. Consequently we do \emph{not} copy published Abd-CT numbers into our
comparison tables. Every baseline we compare against, including FoB, is
re-evaluated by us under this identical pipeline from the authors' released
checkpoints, and figures quoted from the literature are marked as such. This
costs us direct numerical comparability with the published tables and we regard
declaring it as preferable to an unstated protocol mismatch.

\subsection{Training Strategy: Parameter-Isolated Fine-Tuning}
\label{subsec:training}

AdaFoB is initialised from the released FoB Abd-CT checkpoint and fine-tuned
with \textbf{all pretrained parameters frozen}: gradients propagate only to the
newly introduced GAP module (\texttt{refine.*} in our implementation) and to any
parameter absent from the pretrained state dictionary. Concretely, of
\needsnum{total parameters} parameters, \needsnum{trainable parameters} remain
trainable.

Three considerations motivate this choice, and we separate them because they
carry different evidential weight.

\textbf{Compute allocation (established).} The reference FoB schedule is
$36{,}000$ iterations at batch size~$1$. Our budget is a free-tier accelerator
with roughly $30$ GPU-hours per week and $9$--$12$ hour session limits. Under
that constraint a full-network schedule of comparable length is not attainable,
whereas restricting optimisation to the prompt-refinement module concentrates
the entire budget on the parameters that implement the contribution. Training
runs for \needsnum{iterations} iterations with Adam at $10^{-4}$, decayed by
$0.95$ every $1{,}000$ iterations, gradient-norm clipping at $5.0$, and
checkpointing every $500$ iterations with optimiser-state resumption to survive
session preemption.

\textbf{Preserving a converged prompt-localisation prior (design rationale).}
The frozen subnetwork is not a generic feature extractor but a trained
background-prompt localiser whose behaviour is the very baseline we intend to
improve upon. Freezing it makes the comparison between FoB and AdaFoB an
ablation of the sampler alone: the two systems are guaranteed to share
identical encoder features, prototypes and refinement dynamics, so any measured
difference is attributable to prompt geometry rather than to divergent
optimisation of shared weights. We consider this a stronger experimental design
than joint fine-tuning even where compute is not binding.

\textbf{Stability under a randomly initialised head (to be substantiated).}
Attaching a randomly initialised module to a converged network can, in
principle, destabilise the pretrained weights early in training. We do not
claim this effect as an empirical finding, because our current evidence does
not isolate it: our earlier low-scoring runs are attributable to a training
script that never loaded the pretrained checkpoint and ran at $2.8\%$ of the
reference schedule, i.e.\ an undertrained randomly initialised network rather
than a degraded pretrained one. To make a defensible claim we run the direct
comparison in Sec.~\ref{subsec:ablations}: identical initialisation and
schedule, with and without freezing, reporting both final metrics and the
training-loss trajectory. \needsrun{frozen vs.\ full fine-tuning, matched
initialisation and schedule}

\subsection{Evaluation Protocol}
\label{subsec:evaluation}

We report the Dice similarity coefficient (DSC) and the $95$th-percentile
Hausdorff distance (HD95). HD95 is computed between mask \emph{surface} voxels
extracted by morphological erosion; predictions that are empty are assigned a
fixed sentinel distance and the empty-prediction rate is reported alongside, so
that a degenerate predictor cannot be flattered by the distance statistic. Note
that FoB reports Dice only; all HD95 figures in this paper are therefore
produced by us for every method under the identical protocol.

Significance is assessed with the Wilcoxon signed-rank test over paired
per-episode scores, and we report mean $\pm$ standard deviation across folds.
We adopt paired per-episode testing because the effect sizes at stake are on
the order of one Dice point, which is comparable to inter-fold variance and is
therefore not resolvable by comparing fold means alone.

\paragraph{Pipeline-validity gate.}
Before any model comparison we run an \emph{oracle-prompt} control: SAM is
prompted with points sampled from the morphological core of the ground-truth
mask together with background points on the ground-truth offset band. This
removes both prompt generators from the loop and upper-bounds what any prompt
generator could achieve on our data. We report this control and treat it as a
release gate: a comparison of prompt generators is only interpretable when the
oracle clears \needsnum{oracle Dice threshold, target $\geq 0.85$}. We report
this explicitly because prompt-based pipelines can fail silently through
image--label misalignment or intensity-windowing errors, in which case both the
proposed method and its baseline degrade together and the resulting comparison,
though internally consistent, is meaningless.

\paragraph{Prompt attribution analysis.}
To localise where performance is gained or lost we additionally report a
factorial prompt-substitution study crossing $\{$oracle, predicted$\}$ positive
prompts with $\{$oracle, predicted$\}$ background prompts. This isolates the
contribution of the background-prompt pathway -- the pathway that FoB
introduces and that AdaFoB modifies -- from that of the foreground pathway, and
we recommend it as a diagnostic for prompt-generation methods generally.

\subsection{Ablations}
\label{subsec:ablations}

We plan the following, ordered by priority so that the study degrades
gracefully under compute exhaustion:
(i) fixed $N_{p}\in\{5,10,20\}$ versus geometry-derived $N_{p}$, including the
equal-mean-budget control;
(ii) uniform versus curvature-aware placement at matched budget;
(iii) fixed versus scale-adaptive offset $r$;
(iv) frozen versus full fine-tuning (Sec.~\ref{subsec:training});
(v) sensitivity to $[N_{\min}, N_{\max}]$ and to $\nu$;
(vi) an efficiency table reporting trainable parameters, prompt-generation
latency and peak memory against FoB.
\needsrun{all rows}

\subsection{Reproducibility}
\label{subsec:repro}

All experiments use fixed seeds; preprocessing, both prompt generators and the
metric implementations are shared by a single module so that training and
evaluation cannot diverge in their definition of an input. Code, configurations
and per-episode result files will be released. \needsnum{repository URL,
anonymised for review}
